\subsection{覆盖}

设 Y 是道路连通、非空的拓扑空间，\((\mathrm{A}_{i})_{i\in\mathrm{I}}\) 是 Y 由非空道路连通子集组成的一个覆盖，以全序集 I 为指标集。令 \(X=\bigsqcup_{i\in I}A_{i}\) 为族 \((\mathrm{A}_{i})\) 的和空间，\(f:X\to Y\) 为由每个 \(A_{i}\) 到 Y 的规范单射组成的族所诱导的映射。假设映射 f 满足性质 (VK)。特别地，这在下列两种情形成立：

\begin{enumerate}
\def\labelenumi{(\roman{enumi})}
\item
  诸集合 \(A_{i}\)（\(i \in I\)）的内部覆盖 Y（参见 IV, 第~402~页，例 2）；
\item
  空间 Y 是可解圈的，诸空间 \(A_{i}\)（\(i \in I\)）亦然，族 \((\mathrm{A}_{i})_{i \in \mathrm{I}}\) 局部有限，诸 \(A_{i}\) 在 Y 中闭，且它们两两的交局部道路连通（参见 IV, 第~399~页，例 1）。
\end{enumerate}

令 J' 为由三元组 \((i, i', V)\) 组成的集合，其中 i 与 \(i'\) 是 I 的元素，V 是 \(A_i \cap A_{i'}\) 的一个道路连通分支。若 \(j = (i, i', V) \in J'\)，则令 \(i_1(j) = i\)，\(i_2(j) = i'\)，并令 \(\overline{j} = (i', i, V)\)。令 J 为 J' 中由满足 \(i < i'\) 的三元组组成的子集。我们称覆盖的框架为这样的箭图 \(\Gamma\)：其顶点集为 I，箭集为 J，源映射与靶映射分别为映射 \(j \mapsto i_1(j)\) 与 \(j \mapsto i_2(j)\)。我们将把与 \(\Gamma\) 相伴的图等同于这样的图 \(\widetilde{\Gamma}\)：其顶点集为 I，箭集为 \(J \cup \overline{J}\)，源映射与靶映射为映射 \(j \mapsto i_{1}(j)\) 与 \(j \mapsto i_{2}(j)\)，其对合为映射 \(j \mapsto \overline{j}\)。

用 \(p_1\) 与 \(p_2\) 表示纤维方块 \(X \times_{Y} X\) 到 X 上的投影；令 \(\textsf{Γ}\) 为群胚态射偶 \((\varpi(p_1), \varpi(p_2))\)（从 \(\varpi(\mathrm{X} \times_{\mathrm{Y}} \mathrm{X})\) 到 \(\varpi(\mathrm{X})\)）的框架。

引理 1. --- 箭图 \(\Gamma\) 可等同于 \(\textsf{Γ}\) 的一个子箭图；箭图 \(\Gamma\) 与 \(\textsf{Γ}\) 都是连通的。

X 的道路连通分支是诸 \(A_{i}\)（\(i \in I\)）。\(X \times_{Y} X\) 的道路连通分支是诸 \((V \times \{i\}) \times_{Y} (V \times \{i'\})\)（\((i, i', V) \in J'\)）。于是，框架 \(\textsf{Γ}\)（属于偶 \((\varpi(p_1), \varpi(p_2))\)）同构于这样的箭图：其顶点集为 I，箭集为 \(J'\)，源映射与靶映射分别为映射 \(j \mapsto i_1(j)\) 与 \(j \mapsto i_2(j)\)。箭图 \(\textsf{Γ}\) 是连通的（IV, 第~406~页）。此外，这一描述把 \(\Gamma\) 等同于 \(\textsf{Γ}\) 的一个子箭图。我们还注意到，对 \(\textsf{Γ}\) 的每条箭 j，或者 \(j \in \mathrm{Fl}(\widetilde{\Gamma})\)，或者存在 \(i \in I\) 使得 \(j = (i, i, A_i)\)。由此推出，映射 \(\pi_0(\Gamma) \to \pi_0(\textsf{Γ})\)（由 \(\Gamma\) 到 \(\textsf{Γ}\) 的单射导出）是双射，从而 \(\textsf{Γ}\) 是连通的。

对 I 的每个元素 \(i\)，选取一点 \(a(i)\)（在 \(\mathrm{A}_i\) 中）。

对每个元素 \(j = (i, i', \mathrm{V})\)（属于 \(\mathrm{J}\)），选取一点 \(b(j)\)（在 \(\mathrm{V}\) 上）、一条道路 \(\mathrm{B}_1(j)\)，它连接 \(b(j)\) 与 \(a(i)\)（在 \(\mathrm{A}_i\) 上），以及一条道路 \(\mathrm{B}_2(j)\)，它连接 \(b(j)\) 与 \(a(i')\)（在 \(\mathrm{A}_{i'}\) 上）。设 \(j = (i, i', \mathrm{V})\) 是 \(\overline{\mathrm{J}}\) 的一个元素；则 \(\overline{j} = (i', i, \mathrm{V})\) 属于 \(\mathrm{J}\)，且令 \(b(\overline{j}) = b(j)\)，\(\mathrm{B}_1(j) = \mathrm{B}_2(\overline{j})\)，\(\mathrm{B}_2(j) = \mathrm{B}_1(\overline{j})\)。对 \(j \in \mathrm{J}' \cup \overline{\mathrm{J}'}\)，道路 \(\overline{\mathrm{B}_1(j)}\) 与 \(\mathrm{B}_2(j)\)（在 \(\mathrm{Y}\) 上）是可拼接的。令

\[
\mathrm{B} (j) = \overline {{\mathrm{B} _ {1} (j)}} * \mathrm{B} _ {2} (j).\tag{3}
\]

这是 Y 中连接 \(a(i_1(j))\) 与 \(a(i_2(j))\) 的一条道路；我们有关系式 \(\mathrm{B}(\overline{j}) = \overline{\mathrm{B}(j)}\)。

对每个 \(j = (i, i', \mathrm{V}) \in \mathrm{J}'\)，用 \(p_{j,1} \colon V \to A_i\) 与 \(p_{j,2} \colon V \to A_{i'}\) 表示规范单射；再用 \(\varphi_j \colon \pi_1(\mathrm{V}, b(j)) \to \pi_1(\mathrm{A}_i, a(i))\) 与 \(\psi_j \colon \pi_1(\mathrm{V}, b(j)) \to \pi_1(\mathrm{A}_{i'}, a(i'))\) 表示由

\[
\varphi_ {j} (v) = [ \mathrm{B} _ {1} (j) ] ^ {- 1} (p _ {j, 1}) _ {*} (v) [ \mathrm{B} _ {1} (j) ]
\]

以及

\[
\psi_ {j} (v) = [ \mathrm{B} _ {2} (j) ] ^ {- 1} (p _ {j, 2}) _ {*} (v) [ \mathrm{B} _ {2} (j) ],
\]

所定义的群同态，其中 \(v \in \pi_1(\mathrm{V}, b(j))\)（参见 IV, 第~407~页）。

固定 I 的一个元素 \(i_0\)，以及箭图 \(\Gamma\) 中的一个子箭图 T，其相伴图 \(\widetilde{T}\) 是图 \(\widetilde{\Gamma}\) 的极大树。

对 \(i \in I\)，设 \((i_0, j_1, i_1, \ldots, j_n, i)\) 是连接 \(i_0\) 与 \(i\) 的唯一无折返道路（在树 \(\widetilde{T}\) 中），并令

\[
\delta (i) = \left[ \mathrm{B} \left(j _ {1}\right) \right] \left[ \mathrm{B} \left(j _ {2}\right) \right] \dots \left[ \mathrm{B} \left(j _ {n}\right) \right];
\]

它是 Y 中连接 \(a(i_0)\) 与 \(a(i)\) 的一条道路的类。用 \(\alpha_i\) 表示由规范单射导出的从 \(\pi_1(\mathrm{A}_i, a(i))\) 到 \(\pi_1(\mathrm{Y}, a(i))\) 的同态，并令 \(\mu_i: \pi_1(\mathrm{A}_i, a(i)) \to \pi_1(\mathrm{Y}, a(i_0))\) 为由下式定义的群同态：

\[
\mu_ {i} (v) = \delta (i) \alpha_ {i} (v) \delta (i) ^ {- 1}.
\]

最后，令 \(\mu\colon\mathrm{F}(\mathrm{J})\to\pi_{1}(\mathrm{Y},a(i_{0}))\) 为唯一的群同态，使得

\[
\mu (j) = \delta (i _ {1} (j)) [ \mathrm{B} (j) ] \delta (i _ {2} (j)) ^ {- 1}
\]

对所有 \(j \in J\) 成立。存在唯一的群同态

\[
\mathsf {M} \colon \big (\underset {i \in \mathrm{I}} {\ast} \pi_ {1} (\mathrm{A} _ {i}, a (i)) \big) \ast \mathrm{F} (\mathrm{J}) \to \pi_ {1} (\mathrm{Y}, a (i _ {0}))
\]

它与 \(\mu_i\) 在每个 \(\pi_1(A_i, a(i))\)（\(i \in I\)）上、与 \(\mu\) 在 \(F(J)\) 上一致。

令 K' 为由四元组 \((i_{1}, i_{2}, i_{3}, \mathrm{U})\) 组成的集合，其中 \(i_{1}, i_{2}, i_{3}\) 是 I 的元素，U 是 \(A_{i_{1}} \cap A_{i_{2}} \cap A_{i_{3}}\) 的一个道路连通分支。对每个元素 \(k = (i_{1}, i_{2}, i_{3}, \mathrm{U})\)（属于 \(K'\)）与每个偶 \((s, t)\)（由 \(\{1, 2, 3\}\) 中的元素组成），令 \(j_{st}(k) = (i_{s}, i_{t}, \mathrm{V})\)，其中 V 是包含 U 的 \(A_{i_{s}} \cap A_{i_{t}}\) 的道路连通分支；这是 \(J'\) 的一个元素。

令 K 为 \(K'\) 中由四元组 \((i_{1}, i_{2}, i_{3}, \mathrm{U})\)（满足 \(i_{1} < i_{2} < i_{3}\)）组成的子集。对 K 的每个元素 \(k = (i_{1}, i_{2}, i_{3}, \mathrm{U})\)，选取 U 的一点 \(c(k)\)，以及道路 \(\mathrm{C}_{12}(k)\)、\(\mathrm{C}_{23}(k)\) 与 \(\mathrm{C}_{13}(k)\)，使得 \(\mathrm{C}_{st}(k)\) 连接 \(c(k)\) 与 \(b(j_{st}(k))\)（在 \(A_{i_{s}} \cap A_{i_{t}}\) 中），这里 \(s, t \in \{1, 2, 3\}\) 满足 s \textless{} t。

然后，对 \(k \in \mathbf{K}\)，令

\[
\mathrm{L} _ {1} (k) = \overline {{\mathrm{B} _ {1} (j _ {1 3} (k))}} * \overline {{\mathrm{C} _ {1 3} (k)}} * \mathrm{C} _ {1 2} (k) * \mathrm{B} _ {1} (j _ {1 2} (k)),\tag{4}
\]

\[
\mathrm{L} _ {2} (k) = \overline {{\mathrm{B} _ {2} (j _ {1 2} (k))}} * \overline {{\mathrm{C} _ {1 2} (k)}} * \mathrm{C} _ {2 3} (k) * \mathrm{B} _ {1} (j _ {2 3} (k)),
\]

\[
\mathrm{L} _ {3} (k) = \overline {{\mathrm{B} _ {2} (j _ {2 3} (k))}} * \overline {{\mathrm{C} _ {2 3} (k)}} * \mathrm{C} _ {1 3} (k) * \mathrm{B} _ {2} (j _ {1 3} (k)).
\]

对 \(s \in \{1,2,3\}\)，用 \(\lambda_s(k)\) 表示 \(\pi_1(\mathrm{A}_{i_s}, a(i_s))\) 中圈 \(\mathrm{L}_s(k)\) 的类。

命题 1. --- 设 Y 是道路连通的拓扑空间，\((\mathrm{A}_{i})_{i\in\mathrm{I}}\) 是 Y 由非空道路连通子集组成、以全序集 I 为指标集的一个覆盖。假设族 \((\mathrm{A}_{i})_{i\in\mathrm{I}}\) 的和空间到 Y 中的规范映射满足性质 (VK)。那么，上面引入的同态 M 是满射的，其核是包含下列元素的最小正规子群：

\[
\begin{array}{l} \text {(R$_ {1}$)} r _ {1} (j) = j \qquad \text {for } j \text { in } \operatorname{Fl}(\mathsf {T}); \\ \text {(R$_ {2}$)} r _ {2} (j, v) = \varphi_ {j} (v) j \psi_ {j} (v) ^ {- 1} j ^ {- 1} \\ \qquad \qquad \qquad \qquad \qquad \qquad \qquad \qquad \qquad \text {for } j = (i,i^{\prime} ,V)\in \mathsf {J} \text { and } v\in\pi_{1} (V,b(j)); \\ \text {(R$_ {3}$)} r _ {3} (k) = \lambda_ {1} (k) j _ {1 2} (k) \lambda_ {2} (k) j _ {2 3} (k) \lambda_ {3} (k) j _ {1 3} (k) ^ {- 1} \\ \qquad \qquad \qquad \qquad \qquad \qquad \qquad \qquad \text {for } k\in \mathsf {K}. \end{array}
\]

设 X 是族 \((\mathrm{A}_{i})_{i\in\mathrm{I}}\) 的和所成的拓扑空间，\(f\colon X\to Y\) 为由每个 \(A_{i}\) 到 Y 的规范单射组成的族所诱导的映射。由假设，映射 f 满足性质 (VK)；因此我们将对它应用第 1 号的定理 1。于是我们沿用该节的记号，并从为映射 f 定义一个范坎彭数据开始。

X 的道路连通分支是诸集合 \(X_{i} = A_{i} \times \{i\}\)（\(i \in I\)）。这样就把 \(\pi_{0}(X)\) 等同于集合 I。对每个 \(i \in I\)，用 \(\mathfrak{a}(i)\) 表示点 \((a(i), i)\)（属于 \(X_{i}\)）。

令 \(Z = X \times_{Y} X\)。\(Z\) 的道路连通分支是诸集合 \(Z_j = (V \times \{i\}) \times_{Y} (V \times \{i'\})\)，其中 \(j = (i, i', V)\) 取遍 \(J'\)。这样就把 \(\pi_0(Z)\) 等同于集合 \(J'\)。令 \(J_0\) 为 \(J\) 中形如 \((i, i, A_i)\) 的元素组成的集合，于是族 \((J_0, J, \overline{J})\) 是 \(J'\) 的一个划分。对 \(i \in I\) 与 \(j = (i, i, A_i) \in J_0\)，令 \(b(j) = (a(i), a(i))\)，并取 \(\beta_1(j)\) 与 \(\beta_2(j)\) 为 \(a(i)\) 处常值道路的类。对 \(j = (i, i', V) \in J \cup \overline{J}\)，用 \(b(j)\) 表示点 \(((b(j), i), (b(j), i'))\)（属于 \(Z_j\)），并用 \(\beta_1(j)\) 与 \(\beta_2(j)\) 表示道路 \(t \mapsto (B_1(j)(t), i)\) 与 \(t \mapsto (B_2(j)(t), i')\)（在 \(X\) 上）的类。

令 \(W = X \times_{Y} X \times_{Y} X\)。W 的道路连通分支是诸集合 \(W_k = (U \times \{i_1\}) \times_{Y} (U \times \{i_2\}) \times_{Y} (U \times \{i_3\})\)，其中 \(k = (i_1, i_2, i_3, U)\) 取遍 \(K'\)。这样就把 \(\pi_0(W)\) 等同于集合 \(K'\)。

用 \(K_{0}\) 表示 \(K'\) 中形如 \(k = (i, i, i, A_{i})\)（\(i \in I\)）的元素组成的集合。对这样的元素 \(k \in K_{0}\)，令 \(\mathsf{c}(k) = (\mathsf{a}(i), \mathsf{a}(i), \mathsf{a}(i))\)，并取 \(\gamma_{st}(k)\) 为 \((\mathsf{a}(i), \mathsf{a}(i))\) 处常值道路的类。

用 \(K_{1}\) 表示 \(K'\) 中形如 \(k = (i_{1}, i_{2}, i_{3}, V)\) 且集合 \(\{i_{1}, i_{2}, i_{3}\}\) 恰有两个元素的元素组成的集合。设 k 是 \(K'\) 中形如 \((i, i, i', V)\) 的元素，使得 \(j = (i, i', V)\) 属于 \(J \cup \overline{J}\)。这时令

\[
\mathsf {c} (k) = ((b (j), i), (b (j), i), (b (j), i ^ {\prime})), \quad \gamma_ {1 2} (k) = (\beta_ {1} (j), \beta_ {1} (j))
\]

并取 \(\gamma_{13}(k)\) 与 \(\gamma_{23}(k)\) 为 \(\mathsf{b}(j)\) 处常值道路的类。类似地定义 \(c(k)\)、\(\gamma_{12}(k)\)、\(\gamma_{13}(k)\)、\(\gamma_{23}(k)\)（对每个元素 \(k\)，属于 \(\mathrm{K}'\)）。

对 K 的每个元素 \(k = (i_{1}, i_{2}, i_{3}, \mathrm{U})\)，令

\[
\mathsf {c} (k) = ((c (k), i _ {1}), (c (k), i _ {2}), (c (k), i _ {3})) ;
\]

对每个偶 \((s,t)\)（由 \(\{1,2,3\}\) 中不同元素组成），取 \(\gamma_{st}(k)\) 为映射 \(x\mapsto((x,s),(x,t))\) 对道路 \(\mathrm{C}_{st}(k)\) 在 \(\mathrm{A}_{i_{s}(k)}\cap\mathrm{A}_{i_{t}(k)}\) 中的类所取的像。

对每个点 \(x = (x_1, x_2, x_3)\)（属于 \(X \times_Y X \times_Y X\)）与每个置换 \(\sigma \in \mathfrak{S}_3\)，令 \(\sigma(x) = (x_{\sigma^{-1}(1)}, x_{\sigma^{-1}(2)}, x_{\sigma^{-1}(3)})\)。这样就定义了群 \(\mathfrak{S}_3\) 在 W 上的一个作用。对每个 \(k = (i_1, i_2, i_3, U) \in K'\) 与每个置换 \(\sigma \in \mathfrak{S}_3\)，类似地令 \(\sigma(k) = (i_{\sigma^{-1}(1)}, i_{\sigma^{-1}(2)}, i_{\sigma^{-1}(3)}, U)\)；我们有

\[
\sigma \left(\mathrm{W} _ {k}\right) = \left(\mathrm{U} \times \left\{i _ {\sigma^ {- 1} (1)} \right\}\right) \times_ {\mathrm{Y}} \left(\mathrm{U} \times \left\{i _ {\sigma^ {- 1} (2)} \right\}\right) \times_ {\mathrm{Y}} \left(\mathrm{U} \times \left\{i _ {\sigma^ {- 1} (3)} \right\}\right) = \mathrm{W} _ {\sigma (k)}.
\]

设 \(k = (i_1, i_2, i_3, \mathrm{U})\) 是 \(\mathrm{K}'\) 的元素，且 \(i_1, i_2, i_3\) 两两不同。存在唯一的置换 \(\sigma \in \mathfrak{S}_3\) 使得 \(i_{\sigma^{-1}(1)} < i_{\sigma^{-1}(2)} < i_{\sigma^{-1}(3)}\)，亦即使得 \(\sigma(k) = (i_{\sigma^{-1}(1)}, i_{\sigma^{-1}(2)}, i_{\sigma^{-1}(3)}, \mathrm{U})\) 属于 \(\mathrm{K}\)。对 \(s \in \{1, 2, 3\}\)，令 \(c_s(k) = c_{\sigma(s)}(\sigma(k))\)，\(c(k) = (c_1(k), c_2(k), c_3(k))\)，于是 \(c(k) = \sigma^{-1}(c(\sigma(k)))\)。对 \((s, t) \in \{(1, 2), (1, 3), (2, 3)\}\)，定义 \(\mathrm{C}_{st}(k) = \mathrm{C}_{\sigma^{-1}(s)\sigma^{-1}(t)}(\sigma(k))\)；这是连接 \(c(k)\) 与 \(b(j_{\sigma(s)\sigma(t)}(\sigma(k))) = b(j_{st}(k))\) 的一条道路。

用 \(g\) 表示从 \(\Gamma\) 到 \(\textsf{Γ}\) 的箭图态射，它把顶点 \(i \in I\)（属于 \(\Gamma\)）对应为顶点 \(X_i = A_i \times \{i\}\)（属于 \(\textsf{Γ}\)），把箭 \(j = (i, i', V) \in J'\)（属于 \(\Gamma\)）对应为箭 \(Z_j = (V \times \{i\}) \times_Y(V \times \{i'\})\)（属于 \(\textsf{Γ}\)）。映射 Som( \(g\) ) 是双射；映射 Fl( \(g\) ) 是单射，且图 \(\Gamma\) 连通（IV, 第~410~页，引理 1），并且子箭图 T 在 \(g\) 下的像是一个子箭图 \(\mathsf{T}\)（属于 \(\textsf{Γ}\)），其相伴图是图 \(\widetilde{\textsf{Γ}}\) 的极大树。

诸点 \(\mathsf{a}(i)\)（\(i \in I\)）、点 \(\mathsf{b}(j)\)（\(j \in J'\)）、诸点 \(\mathsf{c}(k)\)（\(k \in K'\)）、道路类 \(\beta_{1}(j)\) 与 \(\beta_{2}(j)\)（\(j \in J'\)）、道路类 \(\gamma_{st}(k)\)（\(k \in \mathrm{K}'\)）、子箭图 \(g(\mathrm{T})\)（属于 \(\textsf{Γ}\)）以及元素 \(i_0\)（属于 I）定义了 \(f\) 的一个范坎彭数据。

用 \(\rho\) 表示唯一的群同态

\[
\rho \colon \big (\underset {i \in \mathrm{I}} {*} \pi_ {1} (\mathrm{X}, \mathsf {a} (i)) \big) * \mathrm{F} (\mathrm{J} ^ {\prime}) \to \big (\underset {i \in \mathrm{I}} {*} \pi_ {1} (\mathrm{A} _ {i}, a (i)) \big) * \mathrm{F} (\mathrm{J})
\]

它在 \(\pi_1(\mathrm{X},\mathsf{a}(i))\) 与 \(\pi_1(\mathrm{A}_i,a(i))\) 之间诱导出由 \(\mathrm{A}_i\times \{i\}\) 与 \(\mathrm{A}_i\) 的等同所导出的同构（对每个 \(i\in \mathrm{I}\)），并且使得

\[
\begin{array}{l l} \rho (j) = 1 & \text {for } j \in \mathrm{J} _ {0} \\ \rho (j) = j, \quad \rho (\overline {{j}}) = j ^ {- 1} & \text {for } j \in \mathrm{J}. \end{array}
\]

设 L 为 IV, 第~408~页定理 1 中所定义的群同态。对 \(j = (i, i, \mathrm{A}_i) \in \mathrm{J}_0\)，我们有 \(\mathsf{L}(j) = 1 = \mathsf{M} \circ \rho(j)\)。设 \(j = (i, i', \mathrm{V})\) 是 J 的一个元素；由定义我们有 \(\mathsf{L}(j) = (\mathsf{M} \circ \rho)(j)\)。最后，若 \(j = (i, i', \mathrm{V})\) 是 \(\overline{\mathbf{J}}\) 的一个元素，则 \(\overline{j} \in \mathbf{J}\)，并且可以验证

\[
\mathsf {L} (j) = \mathsf {L} (\overline {{j}}) ^ {- 1} = \mathsf {M} (\rho (\overline {{j}})) ^ {- 1} = \mathsf {M} (\rho (j)).
\]

于是，我们有 M ∘ ρ = L。

同态 L 是满射的（出处同上），因此同态 M 也是满射的。由于同态 \(\rho\) 是满射的，M 的核是 \(\left(\ast\pi_{1}(\mathrm{A}_{i},a(i))\right)*\mathrm{F}(\mathrm{J})\) 中包含在 \(\rho\) 下定理 1（IV, 第~408~页）的关系 \((\mathrm{R}_{1})\)、\((\mathrm{R}_{2})\)、\((\mathrm{R}_{3})\) 所定义元素的像的最小正规子群。只需验证这些像，除了命题 1 的关系 \((\mathrm{R}_{1})\)、\((\mathrm{R}_{2})\)、\((\mathrm{R}_{3})\) 所定义的元素之外，还包括与它们共轭的元素、或与其逆共轭的元素，以及单位元，证明即告完成。

元素 R \(_{1}\) --- 定向树 T 的一条箭形如 Z \(_{j}\)，其中 \(j = (i, i', V) \in J\)；它的像是 F(J) 的元素 \(j\)。

元素 R₂. --- 设 \(j = (i, i', \mathrm{V}) \in \mathrm{J}'\)。若 \(i = i'\)，我们有 \(\rho(\mathsf{r}_2(j, v)) = 1\)（对所有 \(v \in \pi_1(\mathrm{A}_i, \mathsf{a}(i))\)）。若 \(j \in J\)，则 \(\mathsf{r}_2(j, v)\) 的像是元素 \(r_2(j, v) = \varphi_j(v) j \psi_j(v)^{-1} j^{-1}\)（对每个 \(v \in \pi_1(\mathrm{Z}, \mathsf{b}(j))\)）。在余下的情形，我们有 \(\overline{j} \in J\)，且等式

\[
\begin{array}{r l} & {\rho (\mathtt {r} _ {2} (j, v)) = \rho (\varphi_ {j} (v) j \psi_ {j} (v) ^ {- 1} j ^ {- 1})} \\ & {\qquad = [ \mathrm{B} _ {1} (j) ] ^ {- 1} v [ \mathrm{B} _ {1} (j) ] \rho (j) [ \mathrm{B} _ {2} (j) ] ^ {- 1} v ^ {- 1} [ \mathrm{B} _ {2} (j) ] \rho (j) ^ {- 1}} \\ & {\qquad = [ \mathrm{B} _ {2} (\bar {j}) ] ^ {- 1} v [ \mathrm{B} _ {2} (\bar {j}) ] \bar {j} ^ {- 1} [ \mathrm{B} _ {1} (\bar {j}) ] ^ {- 1} v ^ {- 1} [ \mathrm{B} _ {1} (\bar {j}) ] \bar {j}} \end{array}
\]

表明 \(\rho(\mathsf{r}_2(j,v))\) 与 \(\rho(\mathsf{r}_2(\overline{j},v^{-1}))\) 共轭。

元素 R₃. — 设 \(k = (i_{1}, i_{2}, i_{3}, \mathrm{U})\) 是 K' 的一个元素。

若 \(k \in \mathrm{K}_0\)，则 \(i_1 = i_2 = i_3\)，\(\lambda_s(k)\) 对所有 \(s \in \{1,2,3\}\) 是平凡道路的类，且 \(j_{st}(k) \in \mathrm{J}_0\) 对每个偶 \((s,t) \in \{(1,2),(1,3),(2,3)\}\) 成立。此时 \(\rho(\mathsf{r}_3(k))\) 是单位元。

设 \(k\in \mathrm{K}_1\)。若 \(i_1 = i_2\)，则 \(j = (i_{1},i_{3},\mathrm{U})\in \mathrm{J}\cup \overline{\mathrm{J}}\)，并且有

\[
\mathsf {r} _ {3} (k) = \beta_ {1} (j) ^ {- 1} j _ {1 2} \beta_ {1} (j) j _ {2 3} \beta_ {2} (j) ^ {- 1} \beta_ {2} (j) j _ {1 3} ^ {- 1}
\]

它在 \(\rho\) 下的像是单位元。其余情形类似处理。

设 \(i_{1}, i_{2}, i_{3}\) 两两不同。若 \(i_{1} < i_{2} < i_{3}\)，则 \(k \in \mathbf{K}\)，且 \(\rho(\mathsf{r}_{3}(k))\) 的像是元素 \(r_{3}(k)\)。

设 \(\sigma \in \mathfrak{S}_3\) 为把 1 映到 2、把 2 映到 3 的置换。我们有

\[
\begin{array}{r l} \lambda_ {1} (\sigma (k)) & = \beta_ {1} (j _ {1 3} (\sigma (k))) ^ {- 1} \cdot p _ {1, *} (\gamma_ {1 3} (\sigma (k))) ^ {- 1} \cdot \\ & \qquad \cdot p _ {1, *} (\gamma_ {1 2} (\sigma (k))) \cdot \beta_ {1} (j _ {1 2} (\sigma (k))) \\ & = \beta_ {1} (j _ {3 2} (k)) ^ {- 1} \cdot p _ {1, *} (\gamma_ {3 2} (k)) ^ {- 1} \cdot p _ {1, *} (\gamma_ {3 1} (k)) \cdot \beta_ {1} (j _ {3 1} (k)) \\ & = \beta_ {2} (j _ {2 3} (k)) ^ {- 1} \cdot p _ {2, *} (\gamma_ {2, 3} (k)) ^ {- 1} \cdot p _ {2, *} (\gamma_ {1 3} (k)) \cdot \beta_ {2} (j _ {1 3} (k)) \\ & = \lambda_ {3} (k). \end{array}
\]

我们类似地验证 \(\lambda_2(\sigma(k)) = \lambda_1(k)\) 与 \(\lambda_3(\sigma(k)) = \lambda_2(k)\)。于是，

\[
\begin{array}{r l} & {\rho (\mathsf {r} _ {3} (\sigma (k))) = \lambda_ {1} (\sigma (k)) \rho (j _ {1 2} (\sigma (k))) \lambda_ {2} (\sigma (k)) \cdot} \\ & {\qquad \cdot \rho (j _ {2 3} (\sigma (k))) \lambda_ {3} (\sigma (k)) \rho (j _ {1 3} (\sigma (k))) ^ {- 1}} \\ & {\qquad = \lambda_ {3} (k) \rho (j _ {3 1} (k)) \lambda_ {1} (k) \rho (j _ {1 2} (k)) \lambda_ {2} (k) \rho (j _ {3 2} (k)) ^ {- 1}} \\ & {\qquad = \lambda_ {3} (k) j _ {1 3} (k) ^ {- 1} \lambda_ {1} (k) j _ {1 2} (k) \lambda_ {2} (k) j _ {2 3} (k),} \end{array}
\]

这证明 \(\rho(\mathsf{r}_3(\sigma(k)))\) 共轭于

\[
\lambda_ {1} (k) j _ {1 2} (k) \lambda_ {2} (k) j _ {2 3} (k) \lambda_ {3} (k) j _ {1 3} (k) ^ {- 1} = \rho (\mathsf {r} _ {3} (k)).
\]

设 \(\tau \in \mathfrak{S}_3\) 为支集为 \(\{1,2\}\) 的对换。我们有

\[
\lambda_ {1} (\tau (k)) = \lambda_ {2} (k) ^ {- 1}, \quad \lambda_ {2} (\tau (k)) = \lambda_ {1} (k) ^ {- 1} \quad \text {and} \quad \lambda_ {3} (\tau (k)) = \lambda_ {3} (k) ^ {- 1}.
\]

等式

\[
\begin{array}{c} \rho (\mathsf {r} _ {3} (\tau (k))) = \lambda_ {2} (k) ^ {- 1} \rho (j _ {2 1} (k)) \lambda_ {1} (k) ^ {- 1} \rho (j _ {1 3} (k)) \lambda_ {3} (k) ^ {- 1} \rho (j _ {2 3} (k)) ^ {- 1} \\ = \big (\rho (j _ {2 3} (k)) \lambda_ {3} (k) \rho (j _ {1 3} (k)) ^ {- 1} \lambda_ {1} (k) \rho (j _ {1 2} (k)) \lambda_ {2} (k) \big) ^ {- 1} \end{array}
\]

表明 \(\rho(\mathsf{r}_3(\tau(k)))\) 共轭于

\[
\lambda_ {1} (k) \rho (j _ {1 2} (k)) ^ {- 1} \lambda_ {2} (k) \rho (j _ {2 3} (k)) \lambda_ {3} (k) \rho (j _ {1 3} (k)) ^ {- 1} = \rho (\mathsf {r} _ {3} (k)).
\]

由于群 \(\mathfrak{S}_3\) 由置换 \(\tau\) 与 \(\sigma\) 生成，由此推出，对每个 \(k \in \mathrm{K}\) 与每个 \(\sigma \in \mathfrak{S}_3\)，\(\rho(\mathsf{r}_3(\sigma(k)))\) 共轭于 \(\rho(\mathsf{r}_3(k))\) 或其逆。

命题 1 于是得证。

推论 1. --- 在命题 1 的假设之下，再设对每个 \(i \in I\)，从 \(\pi_1(A_i, a(i))\) 到 \(\pi_1(Y, a(i))\) 的同态（由 \(A_i\) 到 Y 的规范单射导出）的像是平凡的。于是由 M 经限制得到的同态 \(M': F(J) \to \pi_1(Y, a(i_0))\) 是满射的，其核是包含元素 \(j \in Fl(T)\) 以及元素 \(j_{12}(k) j_{23}(k) j_{13}(k)^{-1}\)（\(k \in K\)）的最小正规子群。

设 \(\pi\colon (*_{i\in\mathrm{I}}\pi_1(\mathrm{A}_i,a(i)))*\mathrm{F}(\mathrm{J})\to\mathrm{F}(\mathrm{J})\) 为在每个 \(\pi_1(\mathrm{A}_i,a(i))\) 上诱导平凡同态、在 \(\mathrm{F}(\mathrm{J})\) 上诱导恒等同态的唯一同态；它是满射的。设 \(i\in\mathrm{I}\)。由 \(\mathsf{M}\) 的定义以及由规范单射导出的、从 \(\pi_1(\mathrm{A}_i,a(i))\) 到 \(\pi_1(\mathrm{Y},a(i))\) 的同态之像平凡这一假设，可以推出：对每个 \(v\in\pi_1(\mathrm{A}_i,a(i))\)，\(\mathsf{M}(v)\) 是 \(\pi_1(\mathrm{Y},a(i_0))\) 的单位元。因此我们有 \(\mathsf{M}=\mathsf{M}'\circ\pi\)。于是同态 \(\mathsf{M}'\) 是满射的，其核是包含在 \(\pi\) 下命题 1 中所定义的元素 \(r_1(j)\)、\(r_2(j,v)\) 与 \(r_3(k)\) 的像的最小正规子群。对 \(j\in\mathrm{Fl}(\mathrm{T})\)，我们有 \(\pi(r_1(j))=j\)。对所有 \(j\in\mathrm{J}\) 与每个 \(v\in\pi_1(\mathrm{V},b(j))\)，我们有 \(\pi(r_2(j,v))=e\)。最后，对每个 \(k=(i_1,i_2,i_3,\mathrm{U})\in\mathrm{K}\) 与每个 \(s\in\{1,2,3\}\)，\(\lambda_s(k)\) 是 \(\mathrm{A}_{i_s}\) 中一个圈的类；因此 \(\pi(\lambda_s(k))=e\)，从而 \(\pi(r_3(k))=j_{12}(k)j_{23}(k)j_{13}(k)^{-1}\)。推论由此得出。

推论 2. --- 在命题 1 的假设之下，再设对每个 \(i \in I\)，群 \(\pi_1(A_i, a(i))\) 都化为单位元，并且对 I 中两两不同元素的每个三元组 \((i_1, i_2, i_3)\)，集合 \(A_{i_1} \cap A_{i_2} \cap A_{i_3}\) 都是空集。于是同态 \(\mathsf{M}''\colon \mathrm{F}(\mathrm{J} - \mathrm{Fl}(\mathrm{T})) \to \pi_1(\mathrm{Y}, a(i_0))\)（由 \(\mathsf{M}\) 经限制得到）是同构。

设 \(\pi'\colon\mathrm{F}(\mathrm{J})\to\mathrm{F}(\mathrm{J}-\mathrm{Fl}(\mathrm{T}))\) 为这样的同态：它把 \([j]\) 映到 \([j]\)（当 \(j\in\mathrm{J}-\mathrm{Fl}(\mathrm{T})\)），把 \([j]\) 映到单位元（当 \(j\in\mathrm{Fl}(\mathrm{T})\)）；它是满射的，并且我们有 \(\mathsf{M}''\circ\pi'= \mathsf{M}'\)，其中 \(\mathsf{M}'\) 是推论 1 中定义的满射同态。由此推出 \(\mathsf{M}''\) 是满射的，其核是 \(\mathrm{F}(\mathrm{J}-\mathrm{Fl}(\mathrm{T}))\) 中包含在 \(\pi'\) 下出处所述元素的像的最小正规子群。而按构造，\(\pi'(j) = e\)（若 \(j \in \mathrm{Fl}(T)\)），且由假设集合 K 是空集。因此同态 \(M''\) 是同构。

例. --- 1) 关于由两个集合组成的覆盖的情形，见第 3 号。

\begin{enumerate}
\def\labelenumi{\arabic{enumi})}
\setcounter{enumi}{1}
\tightlist
\item
  设 G 是一个图（II, 第~155~页，定义 1）；用 S 表示 G 的顶点集，A 表示其定向边集，o 与 t 表示从 A 到 S 的起点映射与终点映射；对每条定向边 \(a \in A\)，用 \(\overline{a}\) 表示相反的定向边。给集合 S 与 A 赋以离散拓扑；令 X 为空间 S 与空间 I × A 的和空间，并令 ∼ 为 X 上使得 \((u, a) \sim (1 - u, \overline{a})\)、\((0, a) \sim o(a)\) 与 \((1, a) \sim t(a)\)（对所有 \(u \in I\) 与每条定向边 \(a \in A\)）的最细等价关系。商空间 \(|G| = X / \sim\) 称为图 G 的几何实现。我们用 p 表示 X 到 \(|G|\) 上的规范投影。
\end{enumerate}

我们来证明 \(|\mathbf{G}|\) 是局部可缩的。设 \(s \in S\)。用 \(\mathrm{X}_s\) 表示 \(\{s\}\) 与诸子集 \([0,1[ \times \{a\}\)（\(a \in o^{-1}(s)\)）以及诸子集 \(]0,1] \times \{a\}\)（\(a \in t^{-1}(s)\)）的并。令 \(\mathrm{U}_s\) 为 \(\mathrm{X}_s\) 在 \(|\mathbf{G}|\) 中的像；它是 \(p(s)\) 在 \(|\mathbf{G}|\) 中的开邻域，因为 \(\mathrm{X}_s\) 是 \(s\) 在 \(X\) 中的饱和开邻域。令 \(f\) 为从 \(\mathrm{X}_s \times \mathbf{I}\) 到 \(\mathrm{X}_s\) 中的、由下列关系式（对 \(u, v \in \mathbf{I}\) 与 \(a \in A\)）定义的映射

\[
\begin{array}{c} f (s, v) = s \\ f ((u, a), v) = \left\{ \begin{array}{l l} ((1 - v) u, a) & \text {if } 0 \leqslant u <   1 \text {and } o (a) = s, \\ (1 - (1 - v) (1 - u), a) & \text {if } 0 <   u \leqslant 1 \text {and } t (a) = s. \end{array} \right. \end{array}
\]

它是连续的且与等价关系 ∼ 相容。因此，通过过渡到商，它定义了一个映射 \(\varphi_{s}\colon U_{s}\times I\to U_{s}\)，由于 I 局部紧（I, 第~19~页，命题 10），该映射是连续的。它是 \(U_{s}\) 到 \(p(s)\) 上的一个强收缩（III, 第~237~页，定义 6）。

此外设 \(x = (\tau, a) \in ]0, 1[ \times \mathrm{A}\)。令 \(\mathrm{X}_x = ]0, 1[ \times \{a, \overline{a}\}\)，并令 \(\mathrm{U}_x\) 为它在 \(|\mathbf{G}|\) 中（在 \(p\) 之下）的像；它是 \(p(x)\) 在 \(|\mathbf{G}|\) 中的邻域，同胚于 \(]0, 1[\)。从 \(\mathrm{X}_x \times \mathbf{I}\) 到 \(\mathrm{X}_x\) 中的、由 \(((u, a), v) \mapsto ((1 - v)u + v\tau, a)\) 与 \(((u, \overline{a}), t) \mapsto ((1 - v)u + v(1 - \tau), \overline{a})\) 给出的映射是连续的。通过过渡到商，它定义了一个映射 \(\varphi_x: \mathrm{U}_x \times \mathbf{I} \to \mathrm{U}_x\)，它是连续的（出处同上），且是到 \(x\) 处的一个强收缩。

\(|\mathbf{G}|\) 的每个点都是某个点 \(s \in \mathbf{S}\) 的像，或者是 \(\mathbf{I} \times \mathbf{A}\) 中形如 \((\tau, a)\)（其中 \(0 < \tau < 1\)）的点的像。由此推出，\(|\mathbf{G}|\) 的每个点都有在该点处可缩的邻域。换言之，\(|\mathbf{G}|\) 是局部可缩的，特别地是单连通的（IV, 第~346~页，命题 5）。

在 S 上选取一个全序，并考虑由 \((\mathrm{U}_{s})_{s\in\mathrm{S}}\) 组成的 \(|G|\) 的开覆盖。设 s 与 \(s'\) 是 S 的不同元素。交 \(U_{s}\cap U_{s'}\) 是诸 \(p([0,1[,a)\) 的并，其中 a 取遍 G 的以 s 与 \(s'\) 为端点的定向边的集合。于是，覆盖 \((\mathrm{U}_{s})_{s\in\mathrm{S}}\) 的神经等同于箭图 G。此外，对每个 \(s\in S\)，\(U_{s}\) 在 s 处可缩，因此 \(\pi_{1}(\mathrm{U}_{s},p(s))\) 化为单位元。于是由 IV, 第~416~页的推论 2 得出，对每个 \(s\in S\)，\(\pi_{1}(|\mathrm{G}|,p(s))\) 是自由群（尼尔森–施赖尔定理）。

推论 3. --- 在命题 1 的假设之下，再设覆盖 (A \(_{i}\) ) \(_{i\in I}\) 的框架是一个定向树。于是由 M 经限制得到的同态 N: * \(_{i\in I} \pi_{1}(A_{i},a(i)) \to \pi_{1}(Y,a(i_{0}))\) 是满射的；其核是 * \(_{i\in I} \pi_{1}(A_{i},a(i))\) 中包含元素 \(\varphi_{j}(v)\psi_{j}(v)^{-1}\)（对每个 \(j = (i_{1},i_{2},V) \in J\) 与每个 \(v \in \pi_{1}(V,b(j))\)）的最小正规子群。

如果此外诸交 \(A_{i} \cap A_{i'}\)（\(i \neq i'\)）的道路连通分支还是道路单连通的，那么同态 N 是同构。

在该推论的假设之下，与箭图 \(\Gamma\) 相伴的图是一棵树，从而 \(T = \Gamma\)。由此推出，群 \(F(J)\) 在同态 \(M\) 下的像化为单位元。

设

\[
\rho \colon \big (\underset {i \in \mathrm{I}} {\ast} \pi_ {1} (\mathrm{A} _ {i}, a (i)) \big) \ast \mathrm{F} (\mathrm{J}) \to \underset {i \in \mathrm{I}} {\ast} \pi_ {1} (\mathrm{A} _ {i}, a (i))
\]

为在 \(\pi_1(A_i, a(i))\) 上诱导恒等同态且其核包含 F(J) 的唯一群同态。我们有 M = N \(\circ\) \(\rho\)。于是同态 N 是满射的，其核是 \(*_{i \in I} \pi_1(A_i, a(i))\) 中包含在 \(\rho\) 下命题 1 的关系 (R \(_1\) )、(R \(_2\) ) 与 (R \(_3\) ) 所定义元素的像的最小正规子群。我们有 \(\rho(j) = 1\)（对所有 \(j \in \text{Fl}(\Gamma)\)）。由于覆盖 (A \(_i\) ) \(_{i \in I}\) 的框架是一棵树，有 A \(_{i_1} \cap\) A \(_{i_2} \cap\) A \(_{i_3} = \varnothing\)（对 I 中不同元素的每个三元组 \(i_1, i_2, i_3\) ）。由此推出，N 的核是包含元素 \(\varphi_j(v) \psi_j(v)^{-1}\)（对所有 \(j = (i_1, i_2, V) \in J\) 与每个 \(v \in \pi_1(V, b(j))\)）的最小正规子群。

例 3（去掉 n 个点的平面）。--- \(R^{2}-\{0\}\) 的基本群同构于 Z，且道路 \(t \mapsto e^{2\pi it}\) 的类是它的一个生成元（IV, 第~347~页，推论）。更一般地，设 n 是自然数，\(A = \{z_{1}, \ldots, z_{n}\}\) 是 \(R^{2}\) 的由 n 个点组成的集合。令 Y 为空间 \(R^{2}-A\)；我们将证明 Y 的基本群同构于 n 个生成元上的自由群 \(F_{n}\)。对每个 i，令 \(z_{i} = (u_{i}, v_{i})\)。在必要时把 \(z_{i}\) 换成 \(f(z_{i})\)（其中 \(f: R^{2} \to R^{2}\) 是形如 \((u, v) \mapsto (u + \alpha v, v)\) 的同胚）之后，可以假设诸 \(z_{i}\) 的横坐标两两不同。假设 \(u_{1} < \cdots < u_{n}\) 也不失一般性。

令 \(V_{1} = ]-\infty, u_{2}[ \times R\)，\(V_{i} = ]u_{i-1}, u_{i+1}[ \times R\)（对 \(2 \leqslant i \leqslant n-1\)），\(V_{n} = ]u_{n-1}, +\infty[ \times R\)。对 \(1 \leqslant i \leqslant n\)，集合 \(U_{i} = V_{i} - \{z_{i}\}\) 在平面中是开集且同胚于 \(R^{2} - \{0\}\)。族 \((U_{i})_{1 \leqslant i \leqslant n}\) 是空间 \(Y = R^{2} - A\) 的一个开覆盖。交 \(U_{i} \cap U_{j}\) 在 \(|i - j| \geqslant 2\) 时是空集，同胚于 \(R^{2}\)（当 \(|i - j| = 1\)）。由推论 3，Y 的基本群同构于自由群 \(F_{n}\)。

设 \(a\) 是 Y 的一点。对每个整数 \(i \in \{1, \ldots, n\}\)，取一严格正实数 \(r_i\)，它严格小于 \(z_i\) 到诸点 \(z_j\)（\(j \neq i\)）的距离。用 \(v_i\) 表示圈 \(t \mapsto z_i + r_i e^{2\pi it}\) 在 Y 的点 \(z_i + r_i\) 处的类。令 \(\theta_i\) 为一条道路 \(\gamma_i\)（它连接点 \(a\) 与点 \(z_i + r_i\)）在 Y 中的类。如果诸道路 \(\gamma_i\) 都是单射的且它们的像仅在 \(a\) 处相交，那么可以证明，从自由群 F( \(t_1, \ldots, t_n\) ) 到 \(\pi_1(Y, a)\)、使得 \(\varphi(t_i) = \theta_i v_i \theta_i^{-1}\)（对所有 \(i \in \{1, \ldots, n\}\)）的唯一同态是群同构。

类似的论证可以证明：对平面的任意闭离散子集 A，\(R^{2}-A\) 的基本群同构于 F(A)（IV, 第~463~页，习题 1）。

推论 4. --- 在命题 1 的假设之下，设存在 Y 的一个非空道路连通子集 A，使得交 \(\mathrm{A}_i \cap \mathrm{A}_{i'}\) 对每个偶 \((i, i')\)（由 I 中不同元素组成）都等于 A。设 a 是 A 的一点。存在唯一的同态 \(\varphi\)，从群族 \((\pi_1(\mathrm{A}_i, a))_{i \in \mathrm{I}}\) 被 \(\pi_1(\mathrm{A}, a)\) 黏合的和到 \(\pi_1(\mathrm{Y}, a)\) 上，它与由 \(\mathrm{A}_i\) 到 Y 的规范单射导出的同态一致（对每个 \(i \in \mathrm{I}\)）。同态 \(\varphi\) 是同构。

特别地，若群 \(\pi_{1}(\mathrm{A},a)\) 化为单位元，则群族 \((\pi_1(\mathrm{A}_i,a))_{i\in \mathrm{I}}\) 的自由积到 \(\pi_1(\mathrm{Y},a)\) 上的规范同态是同构。

对 \(i \in \mathrm{I}\)，用 \(g_i \colon \pi_1(\mathrm{A}, a) \to \pi_1(\mathrm{A}_i, a)\) 与 \(f_i \colon \pi_1(\mathrm{A}_i, a) \to \pi_1(\mathrm{Y}, a)\) 表示由 \(\mathrm{A}\) 到 \(\mathrm{A}_i\) 以及 \(\mathrm{A}_i\) 到 \(\mathrm{Y}\) 的包含关系导出的规范同态。用 \(*_{\mathrm{A}} \pi_1(\mathrm{A}_i, a)\) 表示诸群 \(\pi_1(\mathrm{A}_i, a)\) 被 \(\pi_1(\mathrm{A}, a)\) 黏合的和。再令 \(p\) 为从 \(*_{i \in \mathrm{I}} \pi_1(\mathrm{A}_i, a)\) 到 \(*_{\mathrm{A}} \pi_1(\mathrm{A}_i, a)\) 上、在 \(\pi_1(\mathrm{A}_i, a)\) 上诱导恒等同态的唯一同态（A, I, 第~80~页，命题 4）。同态 \(p\) 是满射的，并且由幺半群 \(*_{\mathrm{A}} \pi_1(\mathrm{A}_i, a)\) 的定义可知，其核是 \(*_{i \in \mathrm{I}} \pi_1(\mathrm{A}_i, a)\) 中包含元素 \(g_i(v)g_{i'}(v)^{-1}\)（对 \(i, i' \in \mathrm{I}\) 与 \(v \in \pi_1(\mathrm{A}, a)\)）的最小正规子群。

诸同态 \(f_i \circ g_i \colon \pi_1(A, a) \to \pi_1(Y, a)\)（\(i \in I\)）相等。于是由幺半群黏合和的泛性质（A, I, 第~80~页，命题 4），存在唯一的群同态 \(\varphi\)（相应地 \(f\)），从 \(*_{A} \pi_1(A_i, a)\)（相应地从 \(*_{i \in I} \pi_1(A_i, a)\)）到 \(\pi_1(Y, a)\) 中，它诱导出同态 \(f_i\)（到 \(\pi_1(A_i, a)\) 上）。我们有 \(f = \varphi \circ p\)。

对 \(i \in I\)，用 \(u_i\) 表示 A 到 \(A_i\) 中的规范单射，\(v_i\) 表示 \(A_i\) 到 Y 中的规范单射。再用 \(w\) 表示 A 到 Y 中的规范单射。对每个 \(i \in I\)，我们有 \(v_i \circ u_i = w\)，因此 \(\pi_1(v_i, a) \circ \pi_1(u_i, a) = \pi_1(w, a)\)。于是由幺半群黏合和的泛性质（A, I, 第~80~页，命题 4），存在唯一的群同态 \(\varphi\)，从 \(*_{i} \pi_1(A_i, a)\) 到 \(\pi_1(Y, a)\) 上，它诱导出同态 \(\pi_1(v_i, a)\)（到 \(\pi_1(A_i, a)\) 上）。需要证明的是 \(\varphi\) 是同构。

集合 J 等同于由元素偶 \((i, i')\)（属于 I）组成的、满足 \(i < i'\) 的集合。集合 K 等同于由元素三元组 \((i_{1}, i_{2}, i_{3})\)（属于 I）组成的、满足 \(i_{1} < i_{2} < i_{3}\) 的集合。我们选取所有基点 \(a(i)\)、\(b(j)\) 与 \(c(k)\) 都等于 a，所有道路 \(\mathrm{B}(j)\) 与 \(\mathrm{C}_{st}(k)\) 都等于以 a 为像的常值道路。我们再固定一点 \(i_{0} \in I\)。

对每个偶 \((i, i')\)（由 I 的不同点组成），框架 \(\Gamma\)（覆盖 \((\mathrm{A}_{i})\) 的）恰有一条以 i 与 \(i'\) 为端点的箭。\(\Gamma\) 的有一条端点等于 \(i_{0}\) 的箭构成一棵极大定向树 T 的箭。

对每个元素 \(k \in \mathrm{K}\)，我们有 \(r_3(k) = j_{12}(k)j_{23}(k)j_{13}(k)^{-1}\)；令 R 为 \(\mathrm{F(J)}\) 的由元素 \(j \in \mathrm{Fl(T)}\) 与元素 \(r_3(k)\)（\(k \in \mathrm{K}\)）生成的正规子群。我们来证明 \(\mathrm{R} = \mathrm{F(J)}\)。只需证明 J 的每个元素 \(j = (i_1, i_2)\) 都属于 R。由假设，当 \(i_1 = i_0\) 或 \(i_2 = i_0\) 时这是真的。设 \(i_0 < i_1\)，并令 \(k = (i_0, i_1, i_2)\)。它是 K 的元素，且 \(j = j_{23}(k) = j\)。此外，\(j_{12}(k)\) 与 \(j_{13}(k)\) 属于 Fl(T)。由此推出 \(j\) 属于 R。\(i_1 < i_0 < i_2\) 或 \(i_2 < i_0\) 的情形类似处理。

于是由命题 1（IV, 第~412~页）得出，同态 \(f\) 是满射的，其核是包含关系元 \(r_2(j,v) = g_i(v)g_{i'}(v)^{-1}\)（对 \(j = (i,i',\mathrm{A}) \in \mathrm{J}\) 与 \(v \in \pi_1(\mathrm{A},a)\)）的最小正规子群。换言之，\(\operatorname{Ker}(f) = \operatorname{Ker}(p)\)。这就蕴涵 \(\varphi\) 是同构，此即所要证明的。

若 \(\pi_{1}(A,a)\) 化为单位元，则 p 是同构，由此得第二个断言。

例 4. — 设 \(((\mathrm{X}_{i}, x_{i}))_{i \in \mathrm{I}}\) 是一族带基点拓扑空间。我们把族 \(((\mathrm{X}_{i}, x_{i}))_{i \in \mathrm{I}}\) 的楔和定义为并记作 \(\bigvee_{i \in \mathrm{I}} (\mathrm{X}_{i}, x_{i})\)：即族 \((\mathrm{X}_{i})_{i \in \mathrm{I}}\) 的不交并关于把所有点 \((x_{i}, i)\)（\(i \in \mathrm{I}\)）彼此等同起来的等价关系所得的商拓扑空间。用 X 表示这个拓扑空间，用 x 表示诸 \(x_{i}\) 的公共像。假设对每个 \(i \in \mathrm{I}\)，点 \(x_{i}\) 在 \(\mathrm{X}_{i}\) 中是闭的，且诸空间 \(\mathrm{X}_{i}\) 是可解圈的。若 I 有限，则由推论 4，规范同态 \(* \pi_{1}(\mathrm{X}_{i}, x_{i}) \to \pi_{1}(\mathrm{X}, x)\) 是同构。

IV, 第~429~页的注 1 与习题 3（IV, 第~463~页）给出了这个同态是同构的较弱条件。但另见习题 4（IV, 第~464~页）。

